\begin{proofbox}
	\textbf{\textcolor{CProof}{思路提示}}
	采用旋转体体积的圆盘法：体积等于截面面积沿旋转轴的定积分。先由椭圆方程用 \(x\) 表示截面半径平方 \(y^{2}\)，再计算积分 \(\int_{-a}^{a} \pi y^{2}\,dx\)。

	\medskip
	\textbf{\textcolor{CProof}{正式推导}}
	\begin{description}[style=nextline, leftmargin=2.8em, labelsep=0.5em, itemsep=0.55em]
		\item[\textbf{步骤一（写出截面面积）：}]
		      由椭圆方程 \(\frac{x^{2}}{a^{2}}+\frac{y^{2}}{b^{2}}=1\) 解出
		      \[
			      y^{2}=b^{2}\Bigl(1-\frac{x^{2}}{a^{2}}\Bigr).
		      \]
		      绕 \(x\) 轴旋转时，垂直于 \(x\) 轴的截面是圆，半径 \(r=|y|\)，故截面面积函数为
		      \[
			      S(x)=\pi y^{2}=\pi b^{2}\Bigl(1-\frac{x^{2}}{a^{2}}\Bigr),\quad x\in[-a,a].
		      \]
		      \par\textbf{理由：} 旋转体圆盘法，截面半径由曲线到旋转轴的距离确定。
		      \begin{center}
			      % fig02_disk_method_fixed.tex
% 用法：在主文件中提前加载 \usepackage{tikz,amsmath}
% 然后 \input{figures/fig02_disk_method_fixed.tex}
% 说明：这是 TikZ 片段，不包含 \documentclass、\begin{document}、\end{document}。

\begin{center}
	\begin{tikzpicture}[
			x=1cm,
			y=1cm,
			>=stealth,
			line cap=round,
			line join=round,
			every node/.style={font=\small},
			title/.style={font=\bfseries\Large,blue!55!black},
			axis/.style={->,line width=0.8pt,black},
			outline/.style={blue!60!black,line width=0.95pt},
			hidden/.style={blue!50!black,dashed,line width=0.65pt},
			diskline/.style={red!70!black,line width=1.0pt},
			note/.style={red!70!black},
			endpoint/.style={font=\small,fill=blue!6,inner xsep=2pt,inner ysep=1pt},
			formulaText/.style={anchor=north west,align=left,font=\small},
			formulaRed/.style={anchor=north west,align=left,text=red!70!black,font=\bfseries\small}
		]

		% ==================================================
		% 参数设置
		% ==================================================
		\def\semiA{3.60}
		\def\semiB{1.45}
		\def\sliceX{1.20}
		\def\sliceThick{0.18}
		\def\sliceDepth{0.35}

		\pgfmathsetmacro{\sliceY}{\semiB*sqrt(1-(\sliceX*\sliceX)/(\semiA*\semiA))}

		% 手动扩大边界，避免底部公式被裁剪
		\path[use as bounding box] (-4.85,-5.75) rectangle (4.90,2.75);

		% ==================================================
		% 标题
		% ==================================================
		\node[title] at (0,2.35) {圆盘法：垂直于 $x$ 轴切成薄圆盘};

		% ==================================================
		% 旋转椭球体外轮廓
		% ==================================================
		\fill[blue!6] (0,0) ellipse [x radius=\semiA, y radius=\semiB];
		\draw[outline] (0,0) ellipse [x radius=\semiA, y radius=\semiB];

		% 中间虚线截面
		\draw[hidden] (0,0) ellipse [x radius=0.36, y radius=\semiB];

		% x 轴
		\draw[axis] (-4.10,0)--(4.35,0);
		\node at (4.52,0.08) {$x$};

		% 端点标注：下移，避免压在椭圆轮廓和 x 轴上
		\node[endpoint] at (-\semiA,-0.42) {$-a$};
		\node[endpoint] at (\semiA,-0.42) {$a$};

		% ==================================================
		% x 处薄圆盘
		% ==================================================
		% 辅助竖虚线
		\draw[gray!70,dashed,line width=0.65pt]
		(\sliceX,-1.86)--(\sliceX,1.88);

		% 薄圆盘前侧
		\fill[red!8] (\sliceX,0) ellipse [x radius=\sliceThick, y radius=\sliceY];
		\draw[diskline] (\sliceX,0) ellipse [x radius=\sliceThick, y radius=\sliceY];

		% 薄圆盘后侧，增强立体感
		\draw[red!55!black,line width=0.8pt]
		(\sliceX+\sliceDepth,0) ellipse [x radius=\sliceThick, y radius=\sliceY];

		\draw[red!55!black,line width=0.8pt]
		(\sliceX,\sliceY)--(\sliceX+\sliceDepth,\sliceY);

		\draw[red!55!black,line width=0.8pt]
		(\sliceX,-\sliceY)--(\sliceX+\sliceDepth,-\sliceY);

		% 半径 r(x)
		\draw[->,note,line width=0.95pt]
		(\sliceX,0)--(\sliceX,\sliceY);

		\node[note,fill=blue!6,inner xsep=2pt,inner ysep=1pt]
		at (0.45,0.78) {$r(x)=|y|$};

		% x 标注
		\node[fill=blue!6,inner xsep=2pt,inner ysep=1pt]
		at (\sliceX,-1.58) {$x$};

		% dx 标注：放在图形下方，但与公式区保持距离
		\draw[<->,note,line width=0.8pt]
		(\sliceX,-1.82)--(\sliceX+\sliceDepth,-1.82);

		\node[note,fill=blue!6,inner xsep=2pt,inner ysep=1pt]
		at (\sliceX+0.5*\sliceDepth,-2.05) {$dx$};

		% ==================================================
		% 公式说明区：整体下移并分层排版
		% ==================================================
		\node[formulaText] at (-3.95,-2.45)
		{截面面积：\quad $\displaystyle S(x)=\pi r(x)^2=\pi y^2$};

		\node[formulaText] at (-3.95,-3.05)
		{由椭圆方程：\quad $\displaystyle y^2=b^2\left(1-\frac{x^2}{a^2}\right)$};

		\node[formulaRed] at (-3.95,-3.78)
		{$\displaystyle
				\begin{aligned}
					V
					 & =\int_{-a}^{a}S(x)\,dx                                  \\[4pt]
					 & =\int_{-a}^{a}\pi b^2\left(1-\frac{x^2}{a^2}\right)\,dx
				\end{aligned}
			$};

	\end{tikzpicture}
\end{center}

			      \vspace{0.4em}
			      {\small 图：圆盘法截面示意图}
		      \end{center}
		\item[\textbf{步骤二（建立积分式）：}]
		      旋转体体积为截面面积沿 \(x\) 轴的积分，并利用对称性：
		      \[
			      V=\int_{-a}^{a} S(x)\,dx = 2\int_{0}^{a} \pi b^{2}\Bigl(1-\frac{x^{2}}{a^{2}}\Bigr)\,dx.
		      \]
		      \par\textbf{理由：} 被积函数为偶函数，积分区间关于原点对称，可简化为正半轴积分的两倍。

		\item[\textbf{步骤三（计算原函数）：}]
		      提出常数，计算不定积分：
		      \[
			      \int \Bigl(1-\frac{x^{2}}{a^{2}}\Bigr)\,dx = x - \frac{x^{3}}{3a^{2}} + C.
		      \]
		      \par\textbf{理由：} 幂函数积分公式 \(\int x^{n}\,dx = \frac{x^{n+1}}{n+1}\)。

		\item[\textbf{步骤四（代入上下限）：}]
		      计算定积分：
		      \[
			      \begin{aligned}
				      \int_{0}^{a} \Bigl(1-\frac{x^{2}}{a^{2}}\Bigr)\,dx
				       & = \Bigl[x - \frac{x^{3}}{3a^{2}}\Bigr]_{0}^{a} \\
				       & = \Bigl(a - \frac{a^{3}}{3a^{2}}\Bigr) - 0     \\
				       & = a - \frac{a}{3} = \frac{2a}{3}.
			      \end{aligned}
		      \]
		      \par\textbf{理由：} 牛顿-莱布尼茨公式，代入上限减下限。

		\item[\textbf{步骤五（得出体积）：}]
		      将积分结果代回体积表达式：
		      \[
			      V = 2\pi b^{2} \cdot \frac{2a}{3} = \frac{4}{3}\pi a b^{2}.
		      \]
		      \par\textbf{理由：} 合并常数因子，即得所需公式。
	\end{description}

	\medskip
	\textbf{\textcolor{CProof}{补充证明一（沿 y 方向累加椭圆截面）}}
	\textbf{核心识别：}这是同一个绕 \(x\) 轴形成的立体，仅仅改用垂直于 \(y\) 轴的平面切片。切片在 \(xz\) 平面内，\textbf{一般是椭圆，不是圆}。
	\begin{description}[style=nextline, leftmargin=2.8em, labelsep=0.5em, itemsep=0.55em]
		\item[\textbf{步骤一（三维椭球方程）：}]
		      绕 \(x\) 轴旋转时，椭圆上一点到旋转轴的距离 \(|y|\) 变为空间中到 \(x\) 轴的距离 \(\sqrt{y^{2}+z^{2}}\)，故所得实体由
		      \[
		      \frac{x^{2}}{a^{2}}+\frac{y^{2}}{b^{2}}+\frac{z^{2}}{b^{2}}\leq 1
		      \]
		      描述，等号对应外表面。
		      \par\textbf{理由：}同一旋转半径的空间点满足 \(y^{2}+z^{2}=r^{2}\)，转一周后填满内部。

		\item[\textbf{步骤二（固定高度 y=c）：}]
		      在 \(-b<c<b\) 固定 \(y=c\)，截面所在平面为 \(y=c\)，截面边界满足
		      \[
		      \frac{x^{2}}{a^{2}}+\frac{z^{2}}{b^{2}}=1-\frac{c^{2}}{b^{2}},
		      \]
		      等价于
		      \[
		      \frac{x^{2}}{a^{2}(1-c^{2}/b^{2})}
		      +\frac{z^{2}}{b^{2}(1-c^{2}/b^{2})}=1.
		      \]
		      \par\textbf{理由：}这是平面 \(xz\) 中的椭圆方程；\(c=\pm b\) 时截面退化成点，不影响积分。

		\item[\textbf{步骤三（求椭圆截面积）：}]
		      从上式分别读出两个半轴
		      \[
		      A(c)=a\sqrt{1-\frac{c^{2}}{b^{2}}},\qquad
		      B(c)=b\sqrt{1-\frac{c^{2}}{b^{2}}}.
		      \]
		      利用椭圆面积公式 \(S=\pi AB\)，可得
		      \[
		      S(c)=\pi ab\left(1-\frac{c^{2}}{b^{2}}\right).
		      \]
		      \par\textbf{理由：}切片的两条半轴同时乘上 \(\sqrt{1-c^{2}/b^{2}}\)，面积因此乘上 \(1-c^{2}/b^{2}\)。

		\item[\textbf{步骤四（沿 y 积分求体积）：}]
		      厚度为 \(\Delta y\) 的薄片体积约为 \(S(y)\Delta y\)，累加并取极限，得到
		      \[
		      \begin{aligned}
		      V&=\int_{-b}^{b}S(y)\,dy
		      =\pi ab\int_{-b}^{b}\left(1-\frac{y^{2}}{b^{2}}\right)\,dy\\
		       &=\pi ab\left[y-\frac{y^{3}}{3b^{2}}\right]_{-b}^{b}\\
		       &=\pi ab\left(2b-\frac{2b}{3}\right)
		       =\boxed{\frac43\pi ab^{2}}.
		      \end{aligned}
		      \]
		      \par\textbf{理由：}虽然切片方向改变，但各薄片仍完整且不重叠地填满同一个椭球，所以总体积不变。
	\end{description}

	\medskip
	\textbf{\textcolor{CProof}{补充证明二（前述方法四：球体沿 x 方向伸缩，并用方程核对）}}
	\textbf{最快思路：}半径为 \(b\) 的球体积是 \(\frac43\pi b^{3}\)；仅将 \(x\) 方向由 \(b\) 伸缩成 \(a\)，总体积乘 \(a/b\)。
	\begin{description}[style=nextline, leftmargin=2.8em, labelsep=0.5em, itemsep=0.55em]
		\item[\textbf{步骤一（从半径 b 的球出发）：}]
		      用 \(u\) 表示变换前的第一个坐标，球面方程和体积为
		      \[
		      \frac{u^{2}}{b^{2}}+\frac{y^{2}}{b^{2}}+\frac{z^{2}}{b^{2}}=1,
		      \qquad V_{\text{球}}=\frac43\pi b^{3}.
		      \]

		\item[\textbf{步骤二（只伸缩第一个坐标）：}]
		      令 \(x=\frac{a}{b}u\)，而 \(y,z\) 保持不变；反代 \(u=\frac{b}{a}x\)，则
		      \[
		      \begin{aligned}
		      1&=\frac{u^{2}}{b^{2}}+\frac{y^{2}}{b^{2}}+\frac{z^{2}}{b^{2}}\\
		       &=\frac{(bx/a)^{2}}{b^{2}}+\frac{y^{2}}{b^{2}}+\frac{z^{2}}{b^{2}}\\
		       &=\frac{x^{2}}{a^{2}}+\frac{y^{2}}{b^{2}}+\frac{z^{2}}{b^{2}}.
		      \end{aligned}
		      \]
		      \par\textbf{理由：}变换后的边界方程正是绕 \(x\) 轴得到的椭球方程，说明这不是凭外形猜测。

		\item[\textbf{步骤三（证明体积为何乘 a/b）：}]
		      把原球沿 \(u\) 方向切成许多薄圆片。映射 \(x=(a/b)u\) 后，原来位于 \(u\) 的薄圆片被移到 \(x\)，\(y,z\) 均未变化，因此\textbf{截面圆面积不变}；而厚度由 \(\Delta u\) 变成 \(\Delta x=(a/b)\Delta u\)。每个薄片的体积乘上 \(a/b\)，整个球的体积也乘上 \(a/b\)。
		      \par\textbf{理由：}体积是截面积乘厚度再累加；\(a/b>0\)，伸长或压缩都成立。

		\item[\textbf{步骤四（直接得出公式）：}]
		      \[
		      V_{\text{椭球}}=\frac{a}{b}V_{\text{球}}
		      =\frac{a}{b}\cdot\frac43\pi b^{3}
		      =\boxed{\frac43\pi ab^{2}}.
		      \]
	\end{description}
	\textbf{延伸记忆：}从单位球 \(u^{2}+v^{2}+w^{2}\leq 1\) 沿三个方向分别按 \(a,b,c\) 倍伸缩，即得到 \(x^{2}/a^{2}+y^{2}/b^{2}+z^{2}/c^{2}\leq 1\)，因而一般椭球体积为 \(V=\frac43\pi abc\)。在本题中第三条半轴也是 \(b\)，代入便得 \(\frac43\pi ab^{2}\)。

	\medskip

	\textbf{\textcolor{CProof}{结论回扣}}
	\textbf{所得体积公式与结论一致。特别地，当 \(a=b=r\) 时，\(y^{2}=r^{2}(1-x^{2}/r^{2})\)，代入可得 \(V=\frac{4}{3}\pi r^{3}\)，退化为球体积公式。}
\end{proofbox}
